\section{Fundamentos}
\label{sec:fundamentos_zk}

El capítulo anterior describió los mecanismos mediante los cuales una red blockchain
garantiza la integridad de su registro, desde el encadenamiento de bloques por hash hasta las
pruebas de inclusión sobre árboles de Merkle y el modelo de confianza de los clientes ligeros
(Sección~\ref{sec:arquitectura_blockchain}). En todos esos casos, quien verifica una afirmación
sobre la cadena necesita disponer de los datos pertinentes y ejecutar por sí mismo el cómputo
correspondiente. Los sistemas de prueba estudiados en este capítulo, conocidos genéricamente
como pruebas de conocimiento cero, ofrecen una alternativa a este modelo, puesto que permiten
convencer a un verificador de que una afirmación es verdadera, o de que una computación se
ejecutó correctamente, sin que este deba repetir dicha computación. Esta última capacidad no se
deriva de la propiedad de conocimiento cero, sino de la sucintez de los sistemas empleados,
distinción que se precisa en las Secciones~\ref{sec:nocion_zk} y~\ref{sec:snarks_starks}.

La presente sección no pretende ofrecer un tratamiento exhaustivo de la teoría de las pruebas
de conocimiento cero, sino establecer el vocabulario mínimo necesario para comprender las
herramientas analizadas en el resto del capítulo. Para un desarrollo completo de los aspectos
teóricos se remite al lector a la obra de Thaler~\cite{thaler2022proofs}, cuya terminología se
adopta en lo que sigue.

% ---------------------------------------------------------------------------
\subsection{Concepto de conocimiento cero}
\label{sec:nocion_zk}

En su acepción habitual, demostrar una afirmación consiste en presentar evidencia que permita
a otra parte comprobar su veracidad. Si alguien afirma conocer un valor $w$ cuyo hash
criptográfico coincide con un valor público $y$, es decir, un valor tal que $H(w) = y$, la
forma más directa de demostrarlo es entregar $w$ para que el receptor calcule $H(w)$ y lo
compare con $y$. Esta demostración resulta convincente, pero el receptor termina conociendo el
valor que se pretendía mantener en secreto. Goldwasser, Micali y Rackoff demostraron que es
posible convencer a otra parte de una afirmación de este tipo sin transferirle ningún
conocimiento adicional a la validez de la afirmación misma, y denominaron a este tipo de
demostraciones \textit{pruebas de conocimiento cero}~\cite{goldwasser1989knowledge}.
Posteriormente, Goldreich, Micali y Wigderson demostraron que, bajo supuestos criptográficos
estándar, toda afirmación cuya veracidad puede comprobarse eficientemente a partir de un
certificado admite una prueba de este tipo~\cite{goldreich1991proofs}.

El ejemplo anterior permite introducir los cinco términos que se emplean en el resto del
capítulo.

\begin{itemize}
	\item El \textbf{probador} (\textit{prover}), denotado $\mathcal{P}$, es la parte que afirma
	la validez de un enunciado y genera la evidencia que la respalda.
	\item El \textbf{verificador} (\textit{verifier}), denotado $\mathcal{V}$, es la parte que
	examina dicha evidencia y decide si la acepta o la rechaza. Se asume que el verificador
	dispone de recursos de cómputo limitados, por lo que su trabajo debe ser eficiente.
	\item El \textbf{enunciado} (\textit{statement}), denotado $x$, es la afirmación que se desea
	demostrar, expresada a partir de datos públicos conocidos por ambas partes. En el ejemplo, el
	enunciado es que existe un valor cuyo hash es $y$.
	\item El \textbf{testigo} (\textit{witness}), denotado $w$, es la información que el probador
	utiliza para demostrar el enunciado y que no necesita entregar al verificador. En el ejemplo,
	el testigo es la preimagen $w$.
	\item La \textbf{prueba} (\textit{proof}), denotada $\pi$, es el objeto que el probador
	genera a partir del enunciado y del testigo, y que el verificador examina para tomar su
	decisión.
\end{itemize}

Una \textbf{relación} $\mathcal{R}$ es un conjunto de pares $(x, w)$ cuya pertenencia puede
comprobarse eficientemente, y un enunciado $x$ es verdadero cuando existe algún testigo $w$ tal
que $(x, w) \in \mathcal{R}$~\cite{goldreich2001foundations}. En el ejemplo, la relación está
formada por los pares $(y, w)$ que satisfacen $H(w) = y$. Las herramientas que se analizan en
las secciones siguientes reflejan esta distinción mediante las nociones de \textit{entradas
públicas}, que forman parte del enunciado, y \textit{entradas privadas}, que forman parte del
testigo.

Conviene aclarar que el conocimiento cero y la privacidad no son necesariamente lo mismo en las
aplicaciones blockchain. Algunas aplicaciones emplean estas pruebas para ocultar información,
como Zerocash, que valida pagos sin revelar su origen, su destino ni su
monto~\cite{bensasson2014zerocash}. Sin embargo, los sistemas de prueba modernos ofrecen además
una propiedad independiente del conocimiento cero, que consiste en la posibilidad de verificar
una computación en un tiempo muy inferior al que requiere ejecutarla. Walfish y Blumberg
describen este objetivo como la verificación de computaciones sin reejecutarlas, de manera que
una parte con recursos limitados puede delegar un cómputo en otra parte no confiable y
comprobar después, a bajo costo, que el resultado recibido es
correcto~\cite{walfish2015verifying}. En muchas aplicaciones blockchain, los datos involucrados
son públicos y el valor de la prueba reside precisamente en esta segunda propiedad. La
literatura y la industria emplean de todos modos la expresión \textit{prueba de conocimiento
cero} para referirse a estos sistemas, convención que se adopta también en el presente trabajo.

Un ejemplo relacionado con el Capítulo 2 permite apreciar esta diferencia. Una prueba de
inclusión en un árbol de Merkle (Sección~\ref{sec:merkle}) puede formularse como un enunciado,
cuyo contenido es que un elemento $t$ pertenece al árbol de raíz $r$, y un testigo, formado por
los hashes hermanos del camino que une a $t$ con $r$. Ninguno de estos datos es secreto, ya que
todos pueden obtenerse de la cadena. El interés de expresar esta verificación como un enunciado
reside en la posibilidad de combinarla con otras condiciones de validez, como la comprobación de
que la raíz $r$ pertenece a un encabezado certificado por los validadores de la cadena, y de
demostrar el conjunto mediante una única prueba cuya verificación resulte económica.

% ---------------------------------------------------------------------------
\subsection{Propiedades fundamentales}
\label{sec:propiedades_zk}

Un sistema de prueba de conocimiento cero se caracteriza por tres propiedades, cada una de las
cuales protege a una de las partes frente a un comportamiento incorrecto de la
otra~\cite{goldwasser1989knowledge, thaler2022proofs}.

La \textbf{completitud} (\textit{completeness}) garantiza que, si el enunciado es verdadero y el
probador sigue el protocolo utilizando un testigo válido, el verificador acepta la prueba. Esta
propiedad protege al probador honesto, puesto que asegura que una demostración correcta no será
rechazada.

La \textbf{solidez} (\textit{soundness}) garantiza que, si el enunciado es falso, ningún
probador logra que el verificador acepte, salvo con una probabilidad \textit{despreciable}, es
decir, tan pequeña que en la práctica puede considerarse nula. Esta propiedad protege al
verificador frente a un probador deshonesto, denotado $\mathcal{P}^{*}$, que intenta hacer pasar
por verdadera una afirmación falsa. Cuando esta garantía depende de la dificultad de algún problema
computacional, como ocurre en todos los sistemas analizados en este capítulo, el sistema se
denomina técnicamente \textit{argumento} en lugar de \textit{prueba}~\cite{brassard1988minimum}.

En ciertos casos la solidez resulta insuficiente. En el ejemplo de la preimagen, prácticamente
todo valor $y$ posee alguna preimagen, de manera que demostrar su existencia carece de interés.
Lo que se pretende demostrar es que el probador \textit{conoce} una. Esta exigencia se denomina
\textbf{solidez de conocimiento} (\textit{knowledge soundness}) y establece que de todo probador
capaz de convencer al verificador podría obtenerse un testigo válido, lo que implica que dicho
probador efectivamente lo conocía~\cite{bellare1993defining}. Los sistemas que cumplen esta
propiedad se denominan \textit{argumentos de conocimiento}.

El \textbf{conocimiento cero} (\textit{zero-knowledge}) garantiza que el verificador no obtiene
de la prueba ninguna información sobre el testigo más allá de la validez del enunciado. Esta
propiedad protege al probador, y su formalización se basa en mostrar que todo lo que el
verificador observa podría haberse generado a partir del enunciado sin conocer el
testigo~\cite{goldwasser1989knowledge}.

La importancia relativa de cada propiedad depende de la aplicación. En un puente entre cadenas,
el verificador es un contrato inteligente que acuña activos cuando acepta una prueba, y el
adversario principal es un probador que intenta obtener esa acuñación a partir de un evento que
nunca ocurrió en la cadena de origen. En este escenario la propiedad crítica es la solidez,
puesto que su violación permitiría acuñar activos sin respaldo. El conocimiento cero, en cambio,
resulta secundario cuando los datos que se prueban son públicos.

% ---------------------------------------------------------------------------
\subsection{Pruebas interactivas y no interactivas}
\label{sec:interactivas}

En el modelo original de Goldwasser, Micali y Rackoff, probador y verificador intercambian
varios mensajes antes de que el verificador tome su decisión~\cite{goldwasser1989knowledge}. El
verificador envía \textit{desafíos} elegidos al azar, y la solidez se apoya en que el probador
no puede anticiparlos. Este tipo de protocolos se denomina \textbf{interactivo} y presenta
limitaciones importantes para su uso en una red blockchain. La interacción requiere que ambas
partes participen sucesivamente en cada ronda, aunque no necesariamente mediante una conexión
simultánea, y debe repetirse con cada verificador. Además, la transcripción resultante no
constituye, en general, una prueba verificable por terceros, ya que estos no pueden comprobar
que los desafíos se eligieron al azar y no en acuerdo con el probador. En un protocolo de
conocimiento cero, de hecho, todo lo que observa el verificador podría haberse generado sin
conocer el testigo (Sección~\ref{sec:propiedades_zk}), aunque la transferibilidad de una
transcripción depende, en última instancia, del protocolo concreto. En el caso de un contrato inteligente, que se ejecuta de manera
determinista como respuesta a cada transacción, un protocolo de varias rondas podría
implementarse mediante una secuencia de transacciones, pero a costa de gas adicional, de
latencia entre rondas, de la necesidad de conservar estado entre transacciones y de la
dificultad de obtener desafíos aleatorios que el probador no pueda manipular, lo que incrementa
además la complejidad de la implementación.

Una prueba \textbf{no interactiva} consiste, en cambio, en un único mensaje $\pi$ que el
probador genera por sí solo y que cualquier parte puede verificar en cualquier momento. Blum,
Feldman y Micali demostraron que este tipo de pruebas es posible cuando ambas partes comparten
ciertos parámetros públicos generados de antemano mediante un procedimiento denominado
\textbf{configuración inicial} (\textit{setup})~\cite{blum1988noninteractive}. Si ese
procedimiento emplea valores secretos que deben destruirse, porque su conocimiento permitiría
generar pruebas falsas, se habla de una \textbf{configuración confiable}
(\textit{trusted setup}). Si solo emplea aleatoriedad pública, la configuración se denomina
\textbf{transparente}. Esta distinción es una de las diferencias principales entre las familias
de sistemas presentadas en la Sección~\ref{sec:snarks_starks}.

La técnica más difundida para eliminar la interacción es la \textbf{transformación de
Fiat-Shamir}~\cite{fiat1987prove}, aplicable a protocolos de \textit{desafíos públicos}, es
decir, aquellos en los que el verificador se limita a enviar valores aleatorios. La
transformación reemplaza cada desafío por el hash de la transcripción previa, que debe incluir
el enunciado público y los mensajes enviados hasta ese momento por el probador. Un probador
deshonesto puede, sin embargo, calcular el hash de distintos mensajes candidatos antes de elegir
cuál enviar, en busca de uno que le proporcione un desafío favorable, estrategia conocida como
ataque de búsqueda (\textit{grinding}). La seguridad de la transformación suele analizarse en el
modelo de oráculo aleatorio, que idealiza la función de hash como una función perfectamente
aleatoria~\cite{bellare1993random}, y en ese modelo el éxito de dicha búsqueda no depende
únicamente del tamaño del espacio de desafíos, sino también de la solidez del protocolo
original, de sus parámetros y del número de evaluaciones de hash que el adversario puede
realizar. Thaler ilustra este punto señalando que, si el protocolo original ofrece solo 60 bits
de seguridad, unas $2^{60}$ evaluaciones de hash bastan para encontrar un mensaje favorable,
cantidad alcanzable con hardware actual, y recomienda configurar el protocolo con un margen
holgadamente superior a 80 bits~\cite[Sección~5.3]{thaler2022proofs}. Por estos motivos, la
transformación no resulta automáticamente segura para cualquier protocolo ni para cualquier
forma de construir la transcripción. Omitir el enunciado
del cálculo del desafío, por ejemplo, ha dado lugar a ataques
documentados~\cite{bernhard2012pitfalls}. Aplicada correctamente, la transformación permite
generar la prueba sin intervención del verificador. Una prueba así obtenida puede incluirse en una transacción y
ser verificada tanto por un contrato inteligente como por cualquier nodo que valide el bloque,
lo que la convierte en un objeto persistente y verificable públicamente.

% ---------------------------------------------------------------------------
\subsection{De una computación a una prueba}
\label{sec:computacion_prueba}

Esta subsección describe en primer lugar el enfoque basado en circuitos aritméticos y
restricciones, que constituye el fundamento de herramientas como Circom y Noir. En este enfoque,
el desarrollador describe la computación mediante un lenguaje, como Circom o Noir, y un
compilador la transforma en la representación matemática sobre la que opera el sistema de
prueba. La construcción de esa representación se denomina
\textbf{aritmetización}~\cite{thaler2022proofs}. Todos los valores que intervienen en esa
representación son elementos de un \textbf{cuerpo finito} $\mathbb{F}_p$, formado por los enteros
$\{0, 1, \dots, p - 1\}$ con la suma y el producto calculados módulo un número primo $p$, de
manera que, a diferencia de lo que ocurre en un procesador convencional, todo resultado se
reduce módulo $p$. Esta descripción no se aplica de igual modo a todos los sistemas de prueba.
Las zkVM, presentadas al final de esta subsección, no requieren que el desarrollador exprese la
computación de esta forma y emplean internamente representaciones distintas, que se analizan en
la Sección~\ref{sec:sp1}.

\paragraph{Circuitos y restricciones.} Un \textbf{circuito aritmético} representa la
computación como una red de compuertas de suma y de producto sobre $\mathbb{F}_p$. Sin embargo,
un sistema de prueba no necesita repetir la computación, sino comprobar que un conjunto de
valores propuesto por el probador es coherente con ella. Por ese motivo, el circuito se expresa
como un \textbf{sistema de restricciones}, es decir, un conjunto de ecuaciones sobre
$\mathbb{F}_p$ que deben cumplirse simultáneamente y cuyas incógnitas incluyen las entradas, las
salidas y todos los valores intermedios. Para demostrar que se conoce un valor $w$ tal que
$w^{3} + w + 5 = 35$, por ejemplo, la computación puede descomponerse en tres restricciones con
dos variables intermedias $v_1$ y $v_2$.
\begin{equation}
	w \cdot w = v_1, \qquad v_1 \cdot w = v_2, \qquad v_2 + w + 5 = 35 .
\end{equation}
El testigo que satisface este sistema no es solamente el valor $w = 3$, sino la asignación
completa $(w, v_1, v_2) = (3, 9, 27)$. En general, el testigo de un circuito está formado por
todos los valores que surgen al ejecutar la computación. El formato de restricciones más
difundido es el sistema de restricciones de rango uno (\textit{Rank-1 Constraint System},
R1CS)~\cite{parno2013pinocchio}, que se retoma en la Sección~\ref{sec:circom}.

Esta forma de representación tiene dos consecuencias prácticas. La primera es que una
restricción omitida no produce un error de ejecución, sino que amplía el conjunto de testigos
aceptados, lo que permite a un probador deshonesto satisfacer el sistema con valores que no
corresponden a la computación prevista. Los circuitos insuficientemente restringidos constituyen,
de hecho, una categoría recurrente entre las 141 vulnerabilidades reales analizadas por Chaliasos
et al.~\cite{chaliasos2024sok}. La segunda consecuencia es que las operaciones que no son nativas
del cuerpo finito resultan costosas. Las funciones de hash SHA-256 y Keccak-256, empleadas en el
consenso y en las estructuras de datos de Ethereum, operan sobre bits, y cada bit debe
representarse como un elemento de $\mathbb{F}_p$ restringido a los valores $0$ y $1$, por lo que
una única evaluación de estas funciones requiere un número de restricciones muy superior al de
una operación aritmética nativa. Los sistemas actuales mitigan este costo mediante tablas de
búsqueda (\textit{lookup tables}), técnica que se retoma en la Sección~\ref{sec:noir}.

\paragraph{Generación y verificación de la prueba.} Una vez definida la representación, el
proceso completo se desarrolla en cuatro etapas, que la Figura~\ref{fig:flujo_prueba} representa
de manera esquemática.

\begin{figure}[htbp]
	\centering
	\begin{tikzpicture}[
		node distance = 0.9cm and 0.7cm,
		etapa/.style = {draw, rounded corners, align=center, minimum height=1.3cm,
			text width=2.4cm, font=\small},
		dato/.style = {align=center, font=\footnotesize},
		flecha/.style = {-{Stealth[length=2.5mm]}, thick}
	]
		\node[etapa] (comp) {Computación\\y entradas};
		\node[etapa, right=of comp] (wit) {Generación\\del testigo};
		\node[etapa, right=of wit] (prove) {Generación\\de la prueba};
		\node[etapa, right=of prove] (ver) {Verificación\\acepta o\\rechaza};
		\node[etapa, above=of prove] (setup) {Configuración\\inicial};

		\draw[flecha] (comp) -- (wit);
		\draw[flecha] (wit) -- node[above, dato] {$w$} (prove);
		\draw[flecha] (prove) -- node[above, dato] {$\pi$} (ver);
		\draw[flecha] (setup) -- node[left, dato] {parámetros} (prove);
		\draw[flecha] (setup.east) -| node[above, dato, pos=0.25] {parámetros} (ver.north);

		\node[draw, dashed, rounded corners, inner sep=6pt, fit=(comp) (wit) (prove),
			label={[font=\footnotesize]below:Probador $\mathcal{P}$ (\textit{off-chain})}] (zonaP) {};
		\node[draw, dashed, rounded corners, inner sep=6pt, fit=(ver),
			label={[font=\footnotesize]below:Verificador $\mathcal{V}$ (\textit{on-chain})}] (zonaV) {};

		\draw[flecha] (comp.south) -- ++(0, -1.4) -| node[below, dato, pos=0.35]
			{enunciado público $x$} ([xshift=7mm]zonaV.east) |- (ver.east);
	\end{tikzpicture}
	\caption{Etapas de generación y verificación de una prueba.}
	\label{fig:flujo_prueba}
\end{figure}

\begin{enumerate}
	\item \textbf{Configuración inicial.} A partir del sistema de restricciones se generan los
	parámetros públicos que emplean el probador y el verificador. Esta etapa se realiza una única
	vez y, en los sistemas transparentes, no requiere valores secretos.
	\item \textbf{Generación del testigo.} El probador ejecuta la computación sobre las entradas
	concretas y registra todos los valores que intervienen en ella.
	\item \textbf{Generación de la prueba.} El probador transforma el testigo en una prueba $\pi$
	mediante operaciones criptográficas. En muchas construcciones prácticas, esta etapa
	suele ser una de las más costosas del proceso, tanto en tiempo como en memoria.
	\item \textbf{Verificación.} El verificador comprueba la prueba $\pi$ a partir del enunciado
	$x$ y de los parámetros públicos, sin acceso al testigo ni a la computación original, y la
	acepta o la rechaza.
\end{enumerate}

El reparto del trabajo entre ambas partes es marcadamente asimétrico. El costo de generar la
prueba crece al menos de manera lineal con el tamaño de la computación y, en la práctica, suele
superar ampliamente el de ejecutarla, mientras que el costo de verificarla es, en los sistemas que se describen en la
Sección~\ref{sec:snarks_starks}, constante o muy pequeño en comparación. Por ese motivo, en las
aplicaciones blockchain la prueba se genera fuera de la cadena (\textit{off-chain}), en equipos
con recursos suficientes, y se verifica dentro de la cadena (\textit{on-chain}), en un contrato
inteligente que recibe la prueba y el enunciado como datos de una transacción. Conviene
subrayar, además, que la garantía obtenida se limita a que el probador conoce un testigo que
satisface las restricciones para el enunciado dado. Su valor depende, por lo tanto, de que las
restricciones expresen fielmente la afirmación que se pretende establecer, y alcanza únicamente
a los valores que forman parte del enunciado, de modo que todo dato que el verificador utilice para tomar su decisión, como
el destinatario o el monto de una transferencia, debe ser una entrada pública o estar vinculado
criptográficamente a una.

\paragraph{Del circuito a la ejecución de programas.} En el enfoque descrito, cada computación
da lugar a un sistema de restricciones específico, cuyo tamaño queda fijado de antemano. Los
ciclos deben desplegarse con una cantidad máxima de iteraciones conocida al construir el
circuito, las bifurcaciones que dependen de los datos obligan a representar ambas ramas, y toda
modificación de la lógica produce un circuito distinto. Circom y Noir, analizados en las
Secciones~\ref{sec:circom} y~\ref{sec:noir}, pertenecen a este enfoque, aunque difieren en el
nivel de abstracción que ofrecen al desarrollador.

Existe una forma conceptualmente distinta de plantear el problema. En lugar de demostrar que un
testigo satisface las restricciones de una computación particular, puede demostrarse que un
programa arbitrario, ejecutado sobre un procesador determinado, produjo cierto resultado. En
este caso las restricciones no describen el programa sino el procesador, puesto que expresan que
cada paso de la ejecución es una transición válida de acuerdo con la instrucción
correspondiente. El testigo pasa a ser la \textbf{traza de ejecución}, es decir, el registro del
estado del procesador en cada paso, y el programa se convierte en un dato más del enunciado, de
manera que un único sistema de restricciones permite probar la ejecución de cualquier programa.
Ben-Sasson et al.\ materializaron esta idea al construir pruebas sucintas para la ejecución de
programas escritos en el lenguaje C~\cite{bensasson2013snarks}.

Una \textbf{máquina virtual de conocimiento cero} (\textit{zero-knowledge virtual machine},
zkVM) aplica este principio a un conjunto de instrucciones concreto. La zkVM ejecuta un programa
escrito en un lenguaje de propósito general y compilado con herramientas convencionales,
registra su traza y genera una prueba de que la ejecución fue correcta y produjo determinadas
salidas públicas. SP1, analizado en la Sección~\ref{sec:sp1}, ejecuta de
este modo programas escritos en Rust y compilados al conjunto de instrucciones
RISC-V~\cite{succinctsp1}. La diferencia entre ambos enfoques no reside en lo que pueden
demostrar, sino en los compromisos que adoptan. En el enfoque basado en circuitos, el
desarrollador razona en términos de restricciones, lo que exige un conocimiento especializado
pero permite ajustar la representación a la computación. En el enfoque basado en una zkVM, el
desarrollador razona en términos de un programa convencional, lo que facilita el desarrollo, la
reutilización de bibliotecas y la depuración, a cambio de un costo adicional por cada
instrucción ejecutada. Este compromiso es el que motiva la comparación entre Noir y SP1 planteada
en la hipótesis del presente trabajo.

% ---------------------------------------------------------------------------
\subsection{SNARKs y STARKs}
\label{sec:snarks_starks}

Las propiedades descritas en la Sección~\ref{sec:propiedades_zk} no se refieren al tamaño de
la prueba ni al costo de verificarla, aspectos determinantes en las aplicaciones blockchain, ya
que la prueba debe transmitirse en una
transacción y su verificación se paga en gas. Un sistema de prueba es \textbf{sucinto} cuando
tanto el tamaño de la prueba como el tiempo de verificación son muy pequeños en comparación con
el tamaño de la computación probada~\cite{thaler2022proofs}. Las dos familias de sistemas
sucintos más utilizadas en la práctica son los SNARKs y los STARKs.

Ambas familias comparten una idea común, que consiste en codificar la computación como
relaciones entre polinomios y comprobarlas en unos pocos puntos elegidos al azar, lo que se
fundamenta en el principio de Schwartz--Zippel~\cite{schwartz1980fast, zippel1979probabilistic}
y reduce drásticamente el trabajo del verificador. Se diferencian principalmente en el mecanismo
con el que el probador se compromete con esos polinomios, basado en curvas elípticas en los
SNARKs más difundidos, como ocurre con el esquema KZG~\cite{kate2010constant}, y en funciones de
hash en los STARKs, mediante árboles de Merkle y el protocolo FRI~\cite{bensasson2018fri}.

\paragraph{SNARKs.} La sigla SNARK (\textit{Succinct Non-interactive ARgument of Knowledge})
describe un argumento sucinto, no interactivo y con solidez de
conocimiento~\cite{bitansky2012extractable}. Un SNARK que además es de conocimiento cero se
denomina zk-SNARK. Las construcciones más difundidas se basan en curvas elípticas y emplean una operación 
criptográfica denominada \textit{emparejamiento bilineal}, definida entre grupos asociados a
dichas curvas. El sistema de Groth, conocido como Groth16, produce pruebas del orden de pocos
cientos de bytes, cuya verificación requiere un número constante de emparejamientos más un
trabajo proporcional al número de entradas públicas, ya que cada una de ellas debe combinarse con
la clave de verificación mediante una operación sobre la curva~\cite{groth2016size, snarkjs}.
Para un número fijo de entradas públicas y un parámetro de seguridad dado, el costo de
verificación no depende, por lo tanto, del tamaño del circuito.
PLONK, otro sistema de esta familia, admite una configuración universal, válida para cualquier
circuito hasta un tamaño máximo, en lugar de la configuración específica para cada circuito que
requiere Groth16~\cite{gabizon2019plonk}. Noir, el lenguaje analizado en
la Sección~\ref{sec:noir}, delega la generación de pruebas en un backend independiente. Su backend
de referencia, Barretenberg, genera pruebas mediante UltraHonk, un sistema de la familia de
PLONK~\cite{noirlang, barretenbergdocs}.

La principal limitación de construcciones como Groth16, o de PLONK instanciado con compromisos
KZG, es que requieren una configuración confiable, cuyos valores secretos permitirían generar
pruebas falsas si no se destruyen. Este requisito depende de la construcción y no es propio de
todos los SNARKs, ya que existen esquemas basados en compromisos transparentes que no necesitan
una ceremonia de este tipo~\cite{thaler2022proofs}. Para reducir esta
confianza se realizan ceremonias con numerosos participantes, cuyo resultado es seguro siempre
que al menos uno de ellos haya destruido su contribución~\cite{bowe2017scalable}. En contrapartida, estos sistemas ofrecen la
verificación más económica disponible en Ethereum, que desde 2017 incorpora contratos
precompilados para las operaciones sobre la curva elíptica \texttt{alt\_bn128}, también
denominada BN254~\cite{eip196, eip197}.

\paragraph{STARKs.} La sigla STARK (\textit{Scalable Transparent ARgument of Knowledge})
designa una familia de sistemas propuesta por Ben-Sasson et al.\ que prescinde de la
configuración confiable~\cite{bensasson2018scalable}. En lugar de curvas elípticas, los STARKs
se basan únicamente en funciones de hash resistentes a colisiones, lo que los hace además
plausiblemente resistentes a ataques mediante computación cuántica. Son
\textit{transparentes}, en el sentido descrito en la Sección~\ref{sec:interactivas}, y
\textit{escalables}, ya que el tiempo de generación de la prueba crece de manera casi lineal con
el tamaño de la computación, lo que resulta conveniente para computaciones extensas como la ejecución de un
programa completo. Su principal desventaja es el tamaño de la prueba, varios órdenes de magnitud
mayor que el de una prueba Groth16, lo que encarece considerablemente su verificación directa en
Ethereum. La Tabla~\ref{tab:snark_stark} resume las diferencias entre ambas familias.

\begin{table}[htbp]
	\centering
	\caption{Diferencias conceptuales entre SNARKs y STARKs.}
	\label{tab:snark_stark}
	\small
	\begin{tabular}{>{\raggedright\arraybackslash}p{3.6cm} >{\raggedright\arraybackslash}p{5.3cm} >{\raggedright\arraybackslash}p{5.3cm}}
		\toprule
		\textbf{Característica} & \textbf{SNARK} & \textbf{STARK} \\
		\midrule
		Tamaño de la prueba & Muy pequeño & Mayor \\
		\addlinespace
		Verificación & Muy eficiente & Eficiente, pero generalmente más costosa \\
		\addlinespace
		Configuración inicial & Algunas construcciones requieren una configuración confiable &
		Transparente \\
		\addlinespace
		Base criptográfica & Habitualmente curvas elípticas y emparejamientos & Funciones de hash
		y FRI \\
		\addlinespace
		Contexto de uso & Circuitos y verificación compacta & Computaciones extensas y sistemas de
		tipo zkVM \\
		\bottomrule
	\end{tabular}
\end{table}

Ambas familias no son excluyentes. Es posible verificar una prueba STARK dentro de un circuito y
generar sobre esa verificación una prueba SNARK, de manera que el resultado combine la
escalabilidad del probador de los STARKs con la verificación económica de los SNARKs. Esta
composición es la que emplean las zkVM modernas para que sus pruebas puedan verificarse en
Ethereum, y su aplicación en SP1 se analiza en la Sección~\ref{sec:sp1}.

\paragraph{Relación con los puentes entre cadenas.} Los conceptos presentados permiten
reformular el problema de los puentes descrito en la Sección~\ref{sec:puentes}. Un contrato
inteligente desplegado en la cadena de destino necesita establecer la validez de una afirmación
sobre otra cadena, por ejemplo, que una transacción de bloqueo fue incluida en un bloque
certificado por los validadores de la cadena de origen. Un cliente ligero
(Sección~\ref{sec:roles_nodos}) establecería esta afirmación verificando la firma que certifica
el encabezado y la prueba de inclusión de la transacción, pero reproducir esas verificaciones
dentro de un contrato inteligente resulta demasiado costoso. Si ambas verificaciones se expresan
como un único enunciado, el trabajo costoso queda a cargo de un probador fuera de la cadena, y el
contrato solo necesita verificar una prueba sucinta. La solidez del sistema garantiza entonces
que ningún probador puede obtener una prueba válida de un enunciado falso respecto de la relación
compilada, de modo que la aceptación del contrato deja de depender de la honestidad de un
intermediario. Esa garantía no agota, sin embargo, la seguridad del puente, que depende también de
que las restricciones expresen fielmente la afirmación prevista, de los supuestos de seguridad y
de la finalidad del consenso de la cadena de origen, de la vinculación criptográfica entre la
prueba y los datos públicos que utiliza el contrato receptor, de la corrección de la lógica de
ese contrato y de su protección contra la repetición de una misma operación. Xie
et al.\ demostraron la viabilidad práctica de este enfoque~\cite{xie2022zkbridge}. El diseño concreto de esta verificación para el puente
desarrollado en el presente trabajo se expone en el Capítulo~\ref{chp:zkbridge}.
